% !TeX root = ../colloquium-2.tex

\phantomsection\label{sec_Proofs05}
\begin{proof}[Доказательство правил дифференцирования. Вывод производных функций $f(x) = \tg x, f(x) = \sh x$.]
    \subsection*{Доказательство правил дифференцирования.}
    \begin{note}[amber]
        Пусть функции $f$ и $g$ дифференцируемы в точке $a$.
        \begin{enumerate}[label=\arabic*)]
        \item
        функция $\alpha f + \beta g$ дифференцируема в точке $a$ при любых $\alpha, \beta \in \mathbb{R}$ и
        \[
        (\alpha f + \beta g)'(a)=\alpha f'(a)+\beta g'(a)
        \]
        (свойство линейности); \\
            \item
         $f\cdot g$ дифференцируема в точке $a$ и
        \[
        (f\cdot g)'(a)=f'(a)g(a)+f(a)g'(a)
        \]
        (правило Лейбница);
            \item
        если $g\ne 0$ в некоторой окрестности точки $a$, то $\frac{f}{g}$ дифференцируема в точке $a$ и
        \[
        \left(\frac{f}{g}\right)'(a)=\frac{f'(a)g(a)-g'(a)f(a)}{g^2(a)}.
        \]
        \end{enumerate}
        \begin{note}[purple]
            Доказательства.
            \begin{enumerate}[label=\arabic*)]
            \item 
            Найдём производную функции $\alpha f+\beta g$ в точке $a$:
            \[
            \lim_{h\to 0}\frac{(\alpha f+\beta g)(a+h)-(\alpha f+\beta g)(a)}{h}
            \]
            \[
            =\lim_{h\to 0}\left(
            \alpha\frac{f(a+h)-f(a)}{h}
            +
            \beta\frac{g(a+h)-g(a)}{h}
            \right)
            =\alpha f'(a)+\beta g'(a).
            \]
            \item 
            Снова вычислим производную по определению:
            \[
            \lim_{h\to 0}\frac{f(a+h)g(a+h)-f(a)g(a)}{h}
            \]
            \[
            =\lim_{h\to 0}\frac{f(a+h)g(a+h)-f(a)g(a+h)+f(a)g(a+h)-f(a)g(a)}{h}
            \]
            \[
            =\lim_{h\to 0}\left(
            g(a+h)\frac{f(a+h)-f(a)}{h}
            +
            f(a)\frac{g(a+h)-g(a)}{h}
            \right)
            \]
            \[
            = g(a)f'(a)+g'(a)f(a).
            \]
            \item 
            По арифметике пределов имеем:
            \[
            \lim_{h\to 0}\frac{\frac{f(a+h)}{g(a+h)}-\frac{f(a)}{g(a)}}{h}
            \]
            \[
            =\lim_{h\to 0}
            \frac{f(a+h)g(a)-f(a)g(a+h)}{h}\cdot\frac{1}{g(a+h)g(a)}
            \]
            \[
            =\frac{1}{g^2(a)}\lim_{h\to 0}
            \left(
            g(a)\frac{f(a+h)-f(a)}{h}
            -
            f(a)\frac{g(a+h)-g(a)}{h}
            \right)
            \]
            \[
            =\frac{f'(a)g(a)-g'(a)f(a)}{g^2(a)}.
            \]
            \end{enumerate}
        \end{note}
    \end{note}

    \subsection*{Вывод производных функций $f(x) = \tg x, f(x) = \sh x$.}
    \begin{note}[amber]
        \begin{note}[purple]
            \begin{enumerate}[label=\arabic*)]
            \item
            \[
            (\tg x)'=\frac{(\sin x)'\cos x-(\cos x)'\sin x}{\cos^2 x}
            =\frac{\cos^2 x+\sin^2 x}{\cos^2 x}
            =\frac{1}{\cos^2 x}
            =1+\tg^2 x.
            \]
            Аналогично,
            \[
            (\ctg x)'=-\frac{1}{\sin^2 x}=-1-\ctg^2 x.
            \]
            \item 
            Рассмотрим полезные функции: гиперболический синус 
            \[
            \sh x=\frac{e^x-e^{-x}}{2}
            \]
            и гиперболический косинус
            \[
            \ch x=\frac{e^x+e^{-x}}{2}.
            \]
            Согласно правилам дифференцирования
            \[
            (\sh x)'=\frac{1}{2}\big((e^x)'-(e^{-x})'\big)
            =\frac{1}{2}\left(e^x-\left(-e^{-x}\right)\right)
            =\frac{e^x+e^{-x}}{2}
            =\ch x.
            \]
            Аналогично,
            \[
            (\ch x)'=\sh x.
            \]
            \end{enumerate}
        \end{note}
    \end{note}
\end{proof}
